\documentclass[12pt]{book}
\usepackage[utf8]{inputenc}
\usepackage[T1]{fontenc}
\usepackage{amsmath,amssymb,amsfonts}
\usepackage{amsthm}

% Standard operators
\DeclareMathOperator{\Spec}{Spec}

\begin{document}

\setcounter{page}{304}
\section{Residual Complexes.}

Throughout this section, \(X\) will denote a locally noetherian prescheme. If \(x\) is a point of \(X\), and if \(I\) is an injective hull of \(k(x)\) over the local ring \(\sheaf{O}_{x}\), we will denote by \(J(x)\) the quasi-coherent \(\sheaf{O}_{X}\)-module, which is the constant sheaf \(I\) on \(\overline{\{x\}}\), and 0 elsewhere (notation of [II B7]).

\textbf{Definition.} A residual complex on \(X\) is a complex \(K^{\cdot}\) of quasi-coherent injective \(\sheaf{O}_{X}\)-modules, bounded below, with coherent cohomology sheaves, and such that there is an isomorphism
\[
\sum_{p \in \mathbb{Z}} K^{p} \cong \sum_{x \in X} J(x)
\]

\textbf{Example.} If \(X\) is a regular prescheme, then the Cousin complex for the structure sheaf \(\sheaf{O}_{X}\) is a residual complex for \(X\) (see example at end of [IV]).

\begin{proposition}
a) If \(K^{\cdot}\) is a residual complex on \(X\), then its image \(Q(K^{\cdot}) \in D_{c}^{+}(X)\) is a pointwise dualizing complex.

b) If \(R^{\cdot} \in D_{c}^{+}(X)\) is a pointwise dualizing complex on \(X\), then \(E(R^{\cdot})\) is a residual complex on \(X\). (Here \(E\) is the notation of [IV 3], with respect to the filtration \(Z^{\cdot}\) associated with \(R^{\cdot}\)

\setcounter{page}{305}
(cf. [V] Remark 5].)

c) If \(X\) admits a residual complex (or pointwise dualizing complex) with bounded cohomology, then there is a functorial isomorphism
\[
\varphi: R^{\cdot} \to QE\left(R^{\cdot}\right)
\]
for pointwise dualizing complexes. Hence the functor
\[
E: \operatorname{Ptwdual}(X) \to \operatorname{Res}(X)
\]
is an equivalence of the category of pointwise dualizing complexes of \(D_{c}^{+}(X)\) and the category of residual complexes (and morphisms of complexes). Its inverse is \(Q\).
\end{proposition}

\begin{proof}
a) The question is local, so we may assume \(X=\Spec A\), where \(A\) is a noetherian local ring. By [V.3.4] we have only to check that there is an integer \(d\) such that
\[
Ext^{i}\left(k, K^{\cdot}\right)=
\begin{cases}
0 & \text{for } i \neq d \\
k & \text{for } i=d
\end{cases}
\]
where \(k\) is the residue field of \(A\). Since \(k=k(x_{0})\), where \(x_{0}\) is the closed point of \(X\), we have
\[
\Hom \left(k, J\left(x_{0}\right)\right)=k
\]
\[
\Hom (k, J(x))=0 \text{ for } x \neq x_{0}.
\]

\setcounter{page}{306}
The result now follows from the definition of a residual complex.

b) By the pointwise convergence of the spectral sequence [IV.I.G] we see that \(H^{i}(R^{\cdot})=H^{i}(E(R^{\cdot}))\) for all \(i\), and hence \(E(R^{\cdot})\) is a complex which is bounded below and has coherent cohomology. By [IV.I.F] the question of whether it is a residual complex is local, so we may assume \(X\) is the spectrum of a local ring, in which case our result is [V].

c) This follows from [V, Remark 5] and [IV.3.].
\end{proof}

\textbf{Remarks.}
1. In particular, if \(X\) admits a dualizing complex (and hence has finite Krull dimension), the functor
\[
E: \operatorname{Dual}(X) \to \operatorname{Res}(X)
\]
is an equivalence of categories, with inverse \(Q\).

2. I expect that the statement c) is false without the boundedness assumption. That is, there may be two non-isomorphic pointwise dualizing complexes \(R^{\cdot}\) and \(R'^{\cdot}\) such that \(E(R^{\cdot})=E(R'^{\cdot})\). In particular, there may not be a uniqueness theorem similar to [V.3.1] for pointwise dualizing complexes.

3. It follows from the proposition that \(X\) admits a residual complex if and only if it admits a pointwise dualizing complex, so the remarks of [V] apply.

\setcounter{page}{307}
\textbf{Exercise.} Show that there is a uniqueness theorem analogous to [V.3.1] for residual complexes, i.e., two residual complexes can differ only by shifting degrees and tensoring with an invertible sheaf. The touchy point is to show that if \(K^{\cdot}\), \(K'^{\cdot}\) are residual complexes, then \(\Hom^{\cdot}(K^{\cdot}, K'^{\cdot})\) is a complex with coherent cohomology:

We now give two technical results which will be used in the following sections.

\begin{lemma}
Let \(K^{\cdot}\) and \(K'^{\cdot}\) be residual complexes on \(X\). Then to give an isomorphism \(\psi: K^{\cdot} \to K'^{\cdot}\) is equivalent to giving, for each \(x \in X\), an isomorphism
\[
\forall x: Q\left(K_{x}^{\cdot}\right) \to Q\left(K_{x}'\right)
\]
in \(D_{c}^{+}(\Spec \sheaf{O}_{x})\), such that whenever \(x \to y\) is a specialization, is obtained from \(y\) by localization.
\end{lemma}

\begin{proof}
Clearly gives rise to a system \((t_{x})\) as described. Conversely, given the isomorphisms \(v_{x}\) we first note that by c ) of the Proposition above, unique isomorphism \(S_x\) comes from a

\setcounter{page}{308}
\[
\psi_{x}: k_{x}^{*} \to k_{x}'
\]
of the actual residual complexes, and these compatible with localization. And since the codimension functions \(d\) and \(d'\) associated to \(K^{\cdot}\) and \(K'^{\cdot}\) are the same. And since for \(x \neq y\) there are no non-zero maps of \(J(x)\) into \(J(y)\)), our system of isomorphisms \(\bar{t}_{x}\) gives rise to, and is determined by, a collection of isomorphisms
\[
\varphi_{x}: I(x) \to I'(x)
\]
for each \(x \in X\), where \(I(x)\) (resp. \(I'(x)\)) is the (unique) isomorphic to \(J(x)\) subsheaf of \(K^{\cdot}\) (resp. \(K'^{\cdot}\)) for an immediate specialization \(x \to y\), \(\varphi_{x}\) and \(\varphi_{y}\) compatible with the boundary maps of the complexes \(K^{\cdot}\) and \(K'^{\cdot}\). But to give such a collection of isomorphisms \((\varphi_{x})\) to give an isomorphism \(\psi: K^{\cdot} \to K'^{\cdot}\), so we are done.
\end{proof}

\begin{lemma}[Glueing Lemma]
Let \((U_{\nu})\) be an open cover of \(X\), and let \(K_{\nu}^{\cdot}\) be a residual complex on \(U_{\nu}\), for each \(\nu\). Suppose furthermore that for each pair of indices \(\mu,\nu\) we are given an isomorphism

\setcounter{page}{309}
\[
a_{\mu \nu}: K_{\nu}^{*} \to K_{\mu}^{*}
\]
of the restrictions of these complexes to \(U_{\mu \nu}=U_{\mu} \cap U_{\nu}\) such
\[
a_{\mu \nu} a_{\nu \lambda} a_{\lambda \mu}=1
\]
on \(U_{\mu \nu \lambda}\). Suppose finally that the lower bound of the complexes \(K_{\nu}^{*}\) can be chosen uniformly for all \(\nu\). Then there exists a unique residual complex \(K^{\cdot}\) on \(X\), together with isomorphisms
\[
\beta_{\nu}:\left.K^{\cdot}\right|_{U_{\nu}} \to K_{\nu}^{*}
\]
for each \(\nu\), which are compatible with \(a_{\mu \nu}\) on \(U_{\mu \nu}\).
\end{lemma}

\begin{proof}
This is the usual situation for glueing, so it is clear that we can glue the complexes \(K_{\nu}^{*}\) into a global complex \(K^{\cdot}\) on \(X\), which is bounded below. It is also clear that \(K^{\cdot}\) is a complex of quasi-coherent injective \(\sheaf{O}_{X}\)-modules [II.7.16], and has coherent cohomology. Finally since the \(J(x)\) are constant sheaves on irreducible subsets of \(X\), we can glue together local isomorphisms to obtain a global isomorphism
\[
\sum_{p \in \mathbb{Z}} K^{p} \cong \sum_{x \in X} J(x)
\]

\setcounter{page}{310}
Hence \(K^{\cdot}\) is a residual complex.
\end{proof}

\textbf{Remark.} This lemma has the interesting consequence that one can glue dualizing complexes given locally in \(D_{c}^{+}(U_{\nu})\) into a global dualizing complex in \(D_{c}^{+}(X)\), whereas one cannot glue arbitrary objects of the derived category given locally (cf. remarks at end of [III]).

\end{document}